% Worked example for Sec. II.A: a grounded tunable transmon (junctions J_1, J_2)
% capacitively coupled through C_c to an LC resonator, which is driven by a
% voltage source through C_d. Branch order (J_1, J_2, C_c, L_r||C_r, C_d).
% Compile standalone with the draw-circuit driver; \input into main.tex via the
% standalone package. The josephson bipole comes from jj_bipoles.tex, which main.tex also loads.
\documentclass[border=8pt]{standalone}

\usepackage{tikz}
\usepackage{circuitikz}

% Shared circuitikz bipole declarations, copied from the draw-circuit skill template.
% Loaded in the preamble of main.tex and of each standalone figure in figs/.
% Boxed Josephson junction: the co-located E_J || C_J branch with zero loop area.
\ctikzset{
    bipoles/josephson/height/.initial = 0.35,
    bipoles/josephson/width/.initial  = 0.35,
}
\makeatletter
\pgfcircdeclarebipolescaled{misc}
{}
{\ctikzvalof{bipoles/josephson/height}}
{josephson}
{\ctikzvalof{bipoles/josephson/height}}
{\ctikzvalof{bipoles/josephson/width}}
{
    \pgf@circ@setlinewidth{bipoles}{\pgfstartlinewidth}
    \pgfscope
        \pgfpathrectanglecorners{\pgfpoint{\pgf@circ@res@left}{\pgf@circ@res@down}}{\pgfpoint{\pgf@circ@res@right}{\pgf@circ@res@up}}
        \pgfsetfillcolor{white}
        \pgfusepath{fill,stroke}
        \pgfpathmoveto{\pgfpoint{\pgf@circ@res@left}{\pgf@circ@res@down}}
        \pgfpathlineto{\pgfpoint{\pgf@circ@res@right}{\pgf@circ@res@up}}
        \pgfpathmoveto{\pgfpoint{\pgf@circ@res@left}{\pgf@circ@res@up}}
        \pgfpathlineto{\pgfpoint{\pgf@circ@res@right}{\pgf@circ@res@down}}
        \pgfusepath{stroke}
    \endpgfscope
}
\def\pgf@circ@josephson@path#1{\pgf@circ@bipole@path{josephson}{#1}}
\tikzset{josephson/.style = {\circuitikzbasekey, /tikz/to path=\pgf@circ@josephson@path, l=#1}}

\makeatother


\begin{document}
\begin{circuitikz}[line width=0.8pt, font=\small]
    \ctikzset{bipoles/thickness=1, capacitors/scale=0.6, inductors/scale=0.7}
    \tikzset{
        nd/.style={circle, fill, inner sep=1.1pt},
        loop/.style={font=\footnotesize, inner sep=1pt},
    }

    % Top rails (node 1 on the left, node 2 on the right) and the ground rail.
    \coordinate (j1t) at (0,1.5);   \coordinate (j1b) at (0,0);
    \coordinate (j2t) at (1.2,1.5); \coordinate (j2b) at (1.2,0);
    \coordinate (lt) at (3.4,1.5);  \coordinate (lb) at (3.4,0);
    \coordinate (ct) at (4.1,1.5);  \coordinate (cb) at (4.1,0);
    \coordinate (vt) at (6.1,1.5);  \coordinate (vb) at (6.1,0);

    \draw (j1t) to[josephson, l_=$J_1$] (j1b);
    \draw (j2t) to[josephson, l=$J_2$] (j2b);
    \draw (j1t) -- (j2t) to[C, l=$C_c$] (lt);
    \draw (lt) to[L, l_=$L_r$] (lb);
    \draw (ct) to[C, l=$C_r$] (cb);
    \draw (lt) -- (ct) to[C, l=$C_d$] (vt);
    \draw (vb) to[sV, l_=$V(t)$] (vt);
    \draw (j1b) -- (vb);
    \draw (2.3,0) node[ground]{};

    \node[nd] at (j2t) {};
    \node[nd] at (lt) {};
    \node[above, inner sep=2pt] at (0.6,1.5) {$\Phi_1$};
    \node[above, inner sep=2pt] at (3.75,1.5) {$\Phi_2$};

    % Loop fluxes, one per face; each face is closed by branch J_1, C_c, or C_d.
    \node[loop] at (0.6,0.45) {$\odot\,\Phi_x$};
    \node[loop] at (2.3,0.45) {$\odot\,\Phi_y$};
    \node[loop] at (5.15,0.45) {$\odot\,\Phi_z$};
\end{circuitikz}
\end{document}
